% ═══════════════════════════════════════════════════════════════
% LH-02b: tener + edad (Alter angeben) - LEHRERHINWEISE
% Abgeleitet aus LE-02b (Score 9.3/10)
% ═══════════════════════════════════════════════════════════════

\documentclass[11pt, a4paper]{article}

% ═══════════════════════════════════════════════════════════════
% SW-MODUS TOGGLE
% ═══════════════════════════════════════════════════════════════
\newif\ifswmodus
\swmodusfalse    % Standard: Farbe

\usepackage[utf8]{inputenc}
\usepackage[T1]{fontenc}
\usepackage[ngerman]{babel}
\usepackage{helvet}
\renewcommand{\familydefault}{\sfdefault}

\usepackage[margin=2cm]{geometry}
\usepackage{enumitem}
\usepackage{tabularx}
\usepackage{booktabs}
\usepackage{longtable}

\usepackage{xcolor}
\usepackage{amssymb}
\usepackage{tcolorbox}
\tcbuselibrary{skins}

\usepackage{fancyhdr}
\usepackage{lastpage}
\usepackage{needspace}

% Farben
\definecolor{spanisch-color}{HTML}{c41e3a}
\definecolor{grau-color}{HTML}{888888}
\definecolor{hellgrau-color}{HTML}{f5f5f5}
\definecolor{basis-color}{HTML}{2a7a4b}
\definecolor{standard-color}{HTML}{2c5aa0}
\definecolor{erweiterung-color}{HTML}{7b1fa2}
\definecolor{loesung-color}{HTML}{c62828}

% SW-Farben
\definecolor{sw-dark}{HTML}{000000}
\definecolor{sw-gray}{HTML}{666666}
\definecolor{sw-white}{HTML}{ffffff}

% Bedingte Farbzuweisung
\ifswmodus
    \colorlet{spanisch}{sw-dark}
    \colorlet{grau}{sw-gray}
    \colorlet{hellgrau}{sw-white}
    \colorlet{basis}{sw-dark}
    \colorlet{standard}{sw-dark}
    \colorlet{erweiterung}{sw-dark}
    \colorlet{loesung}{sw-dark}
\else
    \colorlet{spanisch}{spanisch-color}
    \colorlet{grau}{grau-color}
    \colorlet{hellgrau}{hellgrau-color}
    \colorlet{basis}{basis-color}
    \colorlet{standard}{standard-color}
    \colorlet{erweiterung}{erweiterung-color}
    \colorlet{loesung}{loesung-color}
\fi

\newcommand{\fachfarbe}{spanisch}

% Variablen
\newcommand{\thema}{tener + edad (Alter angeben)}
\newcommand{\sequenz}{02}
\newcommand{\block}{b}
\newcommand{\fach}{Spanisch}
\newcommand{\klasse}{7}
\newcommand{\dauer}{90 Min. (Doppelstunde)}
\newcommand{\lerngruppengroesse}{24}
\newcommand{\datum}{\today}

% Kopf-/Fußzeile
\pagestyle{fancy}
\fancyhf{}
\renewcommand{\headrulewidth}{0.4pt}
\renewcommand{\footrulewidth}{0.4pt}
\lhead{\fach{} -- Klasse \klasse}
\chead{\textbf{LH-\sequenz\block{} (Lehrerhinweise)}}
\rhead{\datum}
\lfoot{}
\rfoot{Seite \thepage\ von \pageref{LastPage}}

% Befehle
\newcommand{\B}{\textcolor{basis}{\textbf{[B]}}}
\newcommand{\Std}{\textcolor{standard}{\textbf{[S]}}}
\newcommand{\E}{\textcolor{erweiterung}{\textbf{[E]}}}
\newcommand{\loesung}[1]{\textcolor{loesung}{\textbf{#1}}}
\newcommand{\es}[1]{\textcolor{\fachfarbe}{\textbf{#1}}}

% Boxen (SW-kompatibel)
\newtcolorbox{kurzplan}{
    colback=hellgrau,
    colframe=\fachfarbe,
    colbacktitle=white,
    coltitle=\fachfarbe,
    boxrule=1pt,
    arc=0pt,
    fonttitle=\bfseries,
    attach boxed title to top left={yshift=-2mm, xshift=5mm},
    boxed title style={boxrule=0pt, colback=white},
    title=Kurzplan
}

\newtcolorbox{differenzierung}{
    colback=white,
    colframe=grau,
    colbacktitle=white,
    coltitle=grau,
    boxrule=0.5pt,
    arc=0pt,
    fonttitle=\bfseries,
    attach boxed title to top left={yshift=-2mm, xshift=5mm},
    boxed title style={boxrule=0pt, colback=white},
    title=Differenzierung
}

\newtcolorbox{fehlerbox}{
    colback=white,
    colframe=loesung,
    colbacktitle=white,
    coltitle=loesung,
    boxrule=0.5pt,
    arc=0pt,
    fonttitle=\bfseries,
    attach boxed title to top left={yshift=-2mm, xshift=5mm},
    boxed title style={boxrule=0pt, colback=white},
    title=Häufige Fehler
}

\newtcolorbox{materialbox}{
    colback=white,
    colframe=\fachfarbe,
    colbacktitle=white,
    coltitle=\fachfarbe,
    boxrule=0.5pt,
    arc=0pt,
    fonttitle=\bfseries,
    attach boxed title to top left={yshift=-2mm, xshift=5mm},
    boxed title style={boxrule=0pt, colback=white},
    title=Material-Checkliste
}

% ═══════════════════════════════════════════════════════════════
\begin{document}

% Kopfbereich
\begin{center}
    {\Large\bfseries\textcolor{\fachfarbe}{Lehrerhinweise: \thema}}\\[0.3em]
    {\normalsize LH-\sequenz\block{} | \dauer}
\end{center}

\vspace{10pt}

% ═══════════════════════════════════════════════════════════════
% 1. KURZPLAN
% ═══════════════════════════════════════════════════════════════

\section*{1. Stundenverlauf (Kurzplan)}

\textbf{1. Stunde (45 Min.) -- Fokus: Zahlen 1-20}

\begin{tabularx}{\textwidth}{|c|l|X|l|}
\hline
\textbf{Zeit} & \textbf{Phase} & \textbf{Inhalt} & \textbf{Material} \\
\hline
0--5 & Einstieg & ``¿Cuántos años tienes?'' mit Gestik, Vermutungen & -- \\
\hline
5--15 & Input & Zahlen 1-10 mit Fingerzeigen, chorisches Sprechen & Tafel \\
\hline
15--25 & Übung 1 & Zahlen 11-20 mit Rhythmus-Klatsch, Besonderheiten & Zahlenplakat \\
\hline
25--35 & Übung 2 & Zahlen-Bingo & Bingo-Karten \\
\hline
35--40 & Sicherung & Schnell-Quiz ``¿Qué número es?'' & Tafel \\
\hline
40--45 & Überleitung & AB Merkkasten Zahlen ausfüllen & AB-02b \\
\hline
\end{tabularx}

\vspace{1em}

\textbf{2. Stunde (45 Min.) -- Fokus: tener + Alter}

\begin{tabularx}{\textwidth}{|c|l|X|l|}
\hline
\textbf{Zeit} & \textbf{Phase} & \textbf{Inhalt} & \textbf{Material} \\
\hline
0--5 & Aktivierung & Schnelles Zahlen-Quiz (Wiederholung) & -- \\
\hline
5--15 & Input & Verb \textit{tener} einführen, Konjugation an Tafel & Tafel, AB \\
\hline
15--25 & Übung 1 & PA: Alter erfragen ``¿Cuántos años tienes?'' & -- \\
\hline
25--35 & Übung 2 & AB Aufgabe 3 + 4: Alter angeben, Familie beschreiben & AB-02b \\
\hline
35--40 & Sicherung & Ergebnisse besprechen, Fehlerkorrektur & Tafel \\
\hline
40--45 & Abschluss & Zusammenfassung, HA: LP-02b, Vokabeln & -- \\
\hline
\end{tabularx}

% ═══════════════════════════════════════════════════════════════
% 2. DIFFERENZIERUNG
% ═══════════════════════════════════════════════════════════════

\section*{2. Differenzierung}

\begin{differenzierung}
\textbf{Legende:} \B{} Basis (BR) | \Std{} Standard (MR) | \E{} Erweiterung (GY)

\vspace{0.5em}

\textbf{WICHTIG:} Diese Marker sind NUR für die Lehrkraft. Auf Schülermaterial erscheinen KEINE Niveaumarkierungen!
\end{differenzierung}

\vspace{1em}

\begin{tabularx}{\textwidth}{|l|X|X|X|}
\hline
& \B{} \textbf{Basis} & \Std{} \textbf{Standard} & \E{} \textbf{Erweiterung} \\
\hline
\textbf{Aufgabe 1} & Lückentext mit Wortliste & Lückentext selbstständig & + Begründung der Wahl \\
\hline
\textbf{Aufgabe 2} & Zuordnung mit Hilfe & Zuordnung selbstständig & + Zusätzliche Zahlen \\
\hline
\textbf{Aufgabe 3} & Vorgegebene Struktur & Mit Stichworten & Freie Formulierung \\
\hline
\textbf{Aufgabe 4} & 2 Sätze, Wortliste & 3 Sätze & 4+ Sätze, Erweiterungen \\
\hline
\end{tabularx}

\vspace{1em}

\textbf{Konkrete Hinweise:}
\begin{itemize}
    \item \B{} SuS mit LRS: Zeitzugabe (+10 Min.), Zahlenplakat als Hilfe
    \item \B{} DaZ-SuS: Zahlen 1-10 vorab deutsch-spanisch besprechen
    \item \E{} Leistungsstarke: Stammbaum der eigenen Familie erstellen, Zahlen bis 100 einführen
\end{itemize}

% ═══════════════════════════════════════════════════════════════
% 3. HÄUFIGE FEHLER
% ═══════════════════════════════════════════════════════════════

\section*{3. Häufige Fehler}

\begin{fehlerbox}
\begin{tabularx}{\textwidth}{|l|X|X|}
\hline
\textbf{Fehler} & \textbf{Beispiel} & \textbf{Reaktion} \\
\hline
Interferenz ser statt tener & *``Yo soy 12 años'' & Explizit: Spanisch = ``Jahre haben'', nicht ``Jahre sein'' \\
\hline
Verwechslung tiene/tienes & *``Él tienes 15 años'' & Kontrastive Übung: ``Du hast'' vs. ``Er hat'' \\
\hline
Zahlen 11-15 & *``dieciuno'' statt \es{once} & Besondere Formen hervorheben, Rhythmus-Lernen \\
\hline
Aussprache dieciséis & *[dieshi'seis] & Silben: die-ci-séis, betonte Silbe markieren \\
\hline
años vergessen & *``Tengo 12'' & Immer: Zahl + años (Jahre) \\
\hline
\end{tabularx}
\end{fehlerbox}

% ═══════════════════════════════════════════════════════════════
% 4. MATERIAL-CHECKLISTE
% ═══════════════════════════════════════════════════════════════

\section*{4. Material-Checkliste}

\begin{materialbox}
\begin{itemize}[label=$\square$]
    \item AB-02b.pdf (\lerngruppengroesse{} Kopien)
    \item ML-02b.pdf (1× für Lehrkraft)
    \item Zahlenplakat 1-20 (A2 oder digital)
    \item Bingo-Karten (\lerngruppengroesse{} Stück, Zahlen 1-20)
    \item Lehrwerk: Qué pasa 1, S. 24-29
    \item Tafelmarker / Kreide (verschiedene Farben für 11-15)
    \item Optional: LP-02b.html (Tablets/PCs)
    \item Optional: Audio ``Cumpleaños feliz''
\end{itemize}
\end{materialbox}

% ═══════════════════════════════════════════════════════════════
% 5. LÖSUNGEN
% ═══════════════════════════════════════════════════════════════

\section*{5. Lösungen AB-02b}

\textbf{Merkkasten: tener-Konjugation}
\begin{itemize}
    \item yo $\rightarrow$ \loesung{tengo}
    \item tú $\rightarrow$ \loesung{tienes}
    \item él/ella $\rightarrow$ \loesung{tiene}
    \item nosotros $\rightarrow$ \loesung{tenemos}
\end{itemize}

\vspace{1em}

\textbf{Aufgabe 1: tener einsetzen} (4 Punkte)

\begin{tabular}{ll}
a) & Yo \loesung{tengo} 12 años. \\
b) & Mi hermana \loesung{tiene} 15 años. \\
c) & ¿Cuántos años \loesung{tienes} tú? \\
d) & Nosotros \loesung{tenemos} un perro. \\
\end{tabular}

\vspace{1em}

\textbf{Aufgabe 2: Zahlen zuordnen} (4 Punkte)

\begin{tabular}{ll}
doce & $\rightarrow$ \loesung{12} \\
diecisiete & $\rightarrow$ \loesung{17} \\
ocho & $\rightarrow$ \loesung{8} \\
quince & $\rightarrow$ \loesung{15} \\
\end{tabular}

\vspace{1em}

\textbf{Aufgabe 3: Alter angeben} (6 Punkte)

Mögliche Antworten:
\begin{itemize}
    \item a) \loesung{Tengo [12/13/14...] años.}
    \item b) \loesung{Mi madre tiene [z.B. 42] años.}
    \item c) \loesung{Mi hermano/hermana tiene [z.B. 10] años.}
\end{itemize}

\textit{Bewertung: Je 2P (1P korrekte tener-Form, 1P Zahl + años)}

\vspace{1em}

\textbf{Aufgabe 4: Familie beschreiben} (6 Punkte)

Bewertungskriterien:
\begin{itemize}
    \item Mindestens 3 Familienmitglieder genannt: 2P
    \item Korrekte tener-Formen: 2P
    \item Korrekte Zahlen + años: 2P
\end{itemize}

Beispiel:
\begin{itemize}
    \item \loesung{Mi padre se llama Carlos. Tiene 45 años.}
    \item \loesung{Mi madre se llama María. Tiene 42 años.}
    \item \loesung{Mi hermana se llama Ana. Tiene 10 años.}
\end{itemize}

% ═══════════════════════════════════════════════════════════════
% 6. HAUSAUFGABE
% ═══════════════════════════════════════════════════════════════

\section*{6. Hausaufgabe}

\begin{itemize}
    \item Lehrwerk S. 26, Übung 4 (tener einsetzen)
    \item Vokabeln lernen: Zahlen 1-20, tener-Formen
    \item Optional: Lernpfad LP-02b.html (digital)
    \item \E{}: Kurzen Text über die Familie (3-4 Sätze mit Altersangaben)
\end{itemize}

% ═══════════════════════════════════════════════════════════════
% 7. REFLEXIONSNOTIZEN
% ═══════════════════════════════════════════════════════════════

\section*{7. Reflexionsnotizen (nach Durchführung)}

\begin{tabularx}{\textwidth}{|l|X|}
\hline
\textbf{Was lief gut?} & \\[2em]
\hline
\textbf{Was lief nicht?} & \\[2em]
\hline
\textbf{Zeitmanagement} & \\[2em]
\hline
\textbf{Für nächstes Mal} & \\[2em]
\hline
\end{tabularx}

\end{document}
